\documentclass[10pt]{article}
\usepackage[margin=20mm]{geometry}
\usepackage{amsmath,amssymb,mathtools,fontspec,unicode-math,xurl,hyperref}
\IfFontExistsTF{DejaVu Serif}{\setmainfont{DejaVu Serif}}{\setmainfont{Latin Modern Roman}}
\setmathfont{latinmodern-math.otf}
\setlength{\parindent}{0pt}\setlength{\parskip}{6pt}\sloppy
\title{Technical review of
Stabilizing Decomposition of Multiparameter Persistence Modules}
\author{AI-assisted technical review}
\date{}
\begin{document}
\maketitle
Abstract: The central construction of common refinements is supported for positive finite pointwise dimension, subject to a local correction in its exchange argument. Four Minor findings concern an auxiliary distance formula, the zero-module boundary of the stability statements, an unnecessarily restrictive condition in that exchange argument, and two overly strong intermediate assertions in the sharpness proof. The latter assertions admit direct repairs preserving the dimension obstruction. No Major counterexample to the positive-dimensional stability theorem or its order of sharpness was established. The separate conjectures remain conjectures.\par
\section*{Source, definitions and scope}
The reviewed source is Håvard Bakke Bjerkevik’s publisher version of record. Its DOI is 10.1007/s10208-025-09695-w (https://doi.org/10.1007/s10208-025-09695-w), published online on 27 January 2025 and paginated in Foundations of Computational Mathematics 26 (2026), 939–998. The acquired PDF has 60 pages; PDF page \(p\) is printed page \(938+p\). This is an online-2025 assignment.\par
The distinction between a subquotient, an erosion neighborhood, and a refinement is essential. With the paper’s shift convention, \[\operatorname{Er}_\varepsilon(M)\cong
 \operatorname{im}\bigl(M(-\varepsilon)\longrightarrow M(\varepsilon)\bigr).\] A representative \(I/K\) belongs to \(\operatorname{EN}_\varepsilon(M)\) when \[\operatorname{im}(M(-\varepsilon)\to M)\subseteq I\subseteq M,
 \qquad K\subseteq I\cap\ker(M\to M(\varepsilon)).\] A refinement imposes this condition on each indecomposable summand separately. Pointwise finite dimension (pfd) permits the decomposition arguments; a finite uniform bound \(r=\operatorname{supdim}M\) is additionally used in the main stability estimate. The report checks these different hypotheses separately.\par
The source definitions, the distance comparisons, the finite-dimensional pruning mechanism, the exchange step, the sharpness construction, and the constrained-invertibility reformulation were examined. Checks of imported decomposition, matching, and density theorems are limited to their stated use. No author code was executed, and no empirical computation is needed for the counterexamples below.\par
\section*{P1. The printed auxiliary \(d_E\) is not a distance (Minor)}
Definition 4.1 requires the triangle inequality. Definition 4.2 (PDF p. 19/printed p. 957) then defines \[M\leq^*N\quad\Longleftrightarrow\quad
 \dim M_p\leq\dim N_p\quad\text{for every }p,\] and \[d_E(M,N)=\inf\{\varepsilon\geq0:
       \operatorname{Er}_\varepsilon(M)\leq^*N,\ 
       \operatorname{Er}_\varepsilon(N)\leq^*M\}.\] This formula fails the stated requirement, already for pfd one-parameter modules.\par
Let \[M=k_{[0,2)},\qquad
 N=k_{[0,1)}\oplus k_{[1,2)},\qquad Q=0.\] The pointwise dimensions of \(M\) and \(N\) agree, including at the common endpoint \(1\), so \(d_E(M,N)=0\). Indeed, the paper itself uses this pair on PDF p. 20. For a half-open interval, \[\operatorname{Er}_\varepsilon(k_{[a,b)})=
 k_{[a+\varepsilon,b-\varepsilon)},\] with an empty interval denoting zero. Thus \[d_E(M,Q)=1,\qquad d_E(N,Q)=\frac12,\] contradicting \(1\leq 0+\frac12\).\par
The issue is the particular pointwise-dimension formula in this definition. It does not refute an external erosion distance defined on other invariants. The paper’s principal distance \(d_{EN}\) is defined by a common erosion-neighborhood representative; its triangle proof does not use a triangle inequality for \(d_E\). A local correction should remove the distance designation from the displayed formula, or replace the invariant/order and establish the corresponding metric properties.\par
\section*{P2. Exclude \(r=0\) from the stated stability factor (Minor)}
The source explicitly allows \(\operatorname{supdim}M\in\{0,1,\ldots,\infty\}\) on PDF p. 36. Theorem 5.4(ii) and Corollaries 5.5–5.6 (PDF p. 38/printed p. 976) assume only \(r=\operatorname{supdim}M<\infty\).\par
Take the two-parameter modules \[M=0,\qquad N=k_{[0,1)^2},\qquad \varepsilon=1,\] with diagonal shifts. Zero maps give an \(\varepsilon\)-interleaving because both \(2\varepsilon\) structure maps vanish. However, \(\operatorname{Pru}_\varepsilon(M)=0\), while \[\operatorname{EN}_0(N)=\{N\}\quad\text{up to isomorphism}.\] Consequently \(0\) is not a \(2r\varepsilon=0\) refinement of the nonzero module \(N\), contrary to Theorem 5.4(ii). The same example gives \[d_I(0,N)=f_R(0,N)=\frac12>0,\] whereas the right side of Corollary 5.6 is zero.\par
This is a boundary correction, not a failure of the positive-dimensional argument. Require \(r\geq1\), or replace the factor \(r\) by \(\max\{1,r\}\). For \(r=0\), an \(\varepsilon\)-interleaving with zero makes \(N\) \(2\varepsilon\)-trivial, so zero is a \(2\varepsilon\)-refinement of \(N\). The latter formulation therefore supplies the missing case directly. Statements with the same zero factor, including the proposed upper bound in Conjecture 4.20, should carry the corresponding qualification; no claim is made here to prove that conjecture for \(r\geq1\).\par
\section*{P3. Relax the ambient condition in the exchange argument (Minor)}
The proof of Theorem 5.4 replaces indecomposable summands of \(\ker\overline{\psi(\varepsilon)}\) by isomorphic copies which lie inside individual summands of \(\overline N\). The definition of a partial decomposition on PDF p. 40/printed p. 978 additionally requires \[X^S\subseteq\ker\overline{\psi(\varepsilon)}.\] The exchange construction sets \(Y^T=\pi_i(Y)\) and proves \(\pi_i(Y)\subseteq\overline Z_i\) on PDF p. 42. This does not imply the additional kernel inclusion. A projection onto a direct summand need not preserve a kernel: the projection of the diagonal line in \(k^2\) onto the first coordinate is the coordinate line itself.\par
The proof can be formulated without that surplus restriction. Require the replacement copies to be submodules of \(\overline N\), retain their placement in \(\overline Z_j\), their isomorphism types, and both decomposition identities in conditions (i)–(iii). Original summands \(X\) still lie in the kernel. The chain upper-bound argument uses pfd and these decomposition identities. The exchange calculation uses the original \(Y\) and original \(X\)-components when it invokes kernel membership; it does not need the replacement copies to remain there. At maximality, Krull–Schmidt cancellation and the quotients inside the individual \(\overline N_j\) give the required refinement exactly as intended.\par
Thus the written induction invariant is stronger than what its construction establishes. Removing this condition is a local proof repair. The elementary projection example diagnoses the unsupported implication; it is not offered as a persistence-module counterexample to Theorem 5.4.\par
\section*{P4. Two local corrections in the sharpness proof (Minor)}
The following intermediate assertions in the proof of Theorem 5.9 are stronger than necessary. The repaired argument still yields the claimed obstruction to replacing the dimension-dependent factor by a universal constant.\par
\subsection*{Use an inclusion in Lemma 5.10}
On PDF p. 44/printed p. 982, the proof uses \[M_{x_c+(\delta,\delta)}=k^{c+n-2}.\] For \(n=3\), \(c=1\), and \(\delta=0\), the actual generators in the preceding page give \(M_{x_1}=\langle e_1\rangle\), while the displayed right side is \(\langle e_1,e_2\rangle\). The valid statement is \[M_{x_c+(\delta,\delta)}\subseteq k^{c+n-2},
 \qquad 0\leq\delta<n-1.\] Indeed, \(\lfloor\delta\rfloor\leq n-2\), which bounds the indices of generators present at this point. The subsequent flag argument needs only this inclusion.\par
\subsection*{Use the rank bound for the second image}
On PDF pp. 45–46, write \(T=M_{0\to1}\). The proof asserts that \(e_{2n}\) never belongs to \(\operatorname{im}(T-f)\). Near the terminal boundary some earlier coordinates have already died, so this is false. For \(n=2\), consider the map from \((24,24)\) to \((25,25)\). Its domain contains all four coordinate vectors, but \(I_1\) has disappeared at the target. Consequently \[(T-f)e_1=-e_2,\qquad
 (T-f)e_2=e_2-e_3,\qquad
 (T-f)e_3=e_3-e_4.\] Their span contains \(e_4\).\par
Nevertheless the required estimate \(\dim\operatorname{im}(T-f)\leq2n-1\) holds. If the domain has dimension below \(2n\), it follows immediately. Otherwise the domain is all of \(k^{2n}\) and \((T-f)e_{2n}=0\), so rank–nullity gives the estimate. The other image bound and the block interleaving of Lemma 5.12 remain valid. These corrections repair the route to Theorem 5.9; they do not provide a counterexample to that theorem.\par
\section*{Positive checks and limits}
\subsection*{Erosion neighborhoods and distances}
Lemma 4.4 follows by taking inverse images of the submodule and kernel defining the second quotient. The resulting numerator contains \(\operatorname{im}(M(-\varepsilon-\delta)\to M)\), and the denominator is killed by the \((\varepsilon+\delta)\) shift.\par
For Lemma 4.5, write two representatives inside the middle module as \(A_i/B_i\). Their common refinement is obtained from \[A_3=A_1\cap A_2,\qquad B_3=(B_1+B_2)\cap A_3.\] Under \(A_1\to A_1/B_1\), the image of \(A_3\) contains the required shifted image, and the image of \(B_3\) lies in the shifted kernel. This proves the triangle inequality for \(d_{EN}\) independently of P1.\par
Theorem 4.7 was checked by using the two maps from \(M(-\varepsilon)\oplus N(-\varepsilon)\) associated with an interleaving. Naturality and the interleaving identities identify the inverse images of the two quotient kernels, yielding the same quotient in both erosion neighborhoods. Conversely, interleavings through a common representative compose. Thus \(d_{EN}\leq d_I\leq2d_{EN}\) is supported.\par
For one-parameter half-open bars, the induced-matching proof of Theorem 4.9 controls both endpoint changes by \(\varepsilon\); unmatched bars have length at most \(2\varepsilon\). Bars of exactly that length need not be matched. The general decorated-endpoint case is delegated by the source to its cited matching theorem. The counterexamples in Lemmas 4.10–4.11, refinement composition, the pruning-distance triangle argument, and the thin-summand argument in Theorem 4.22 were also checked. In that last argument, the union of two pointwise subspaces is a submodule because each ambient space has dimension at most one.\par
\subsection*{Pruning, stability and sharpness}
For \(1\leq r<\infty\), the descending and ascending pointwise chains in Lemma 5.3 stabilize after at most \(r\) strict changes. Their image/kernel inclusions, the induced quotient maps, and the splitting step support Theorem 5.4, with P3’s repaired exchange invariant. This is a finite-dimensional argument; the general-module countable construction in Lemma 5.2 was not independently certified without a finiteness hypothesis. Existence of a largest admissible numerator and smallest admissible denominator can instead be obtained by sums and intersections of admissible submodules.\par
The flag constraints in Lemma 5.10 and the block-matrix identities in Lemma 5.12 support the sharpness mechanism after P4’s corrections. No numerical sampling is used as a substitute for those arguments.\par
\subsection*{Constrained invertibility and conjectural statements}
The upset-module reduction and finite constrained-invertibility construction in Section 6 were checked. Theorem 6.14 moves from pfd upset decompositions to finite matrix instances without spelling out the reduction. A finite-witness argument fills this step: interleavings give mutually inverse column-finite matrices; protect the finitely many vertices reachable within the relevant bounded path length from a finite Hall-test set, and extend the maps on these protected columns to inverse finite matrices. No new short neighbors of the test set appear. The finite assumption then implies Hall’s condition. The matching graph is locally finite by pfd, so the infinite matching step applies. This supplies counterevidence to treating the omitted reduction as a theorem failure.\par
The paper does not claim to prove Conjectures 4.20, 6.1 or 6.10. The review does not resolve them, establish every imported result, or certify a polynomial-time implementation of pruning. Those limits are separate from the four local findings.\par
\section*{Supporting checks}
The appropriate technical assessment is Minor: the auxiliary distance formula and zero-module scope need correction, and two proof passages need the explicit local repairs above. The positive-dimensional core and the dimension-dependent obstruction retain substantive support. This assessment is an initial scientific review, with independent acceptance still pending.\par
No usable substantive predecessor report for this paper was found in the accessible preserved inventory. The acquired source PDF, manifest, and historical extraction were reused with their original identities. The present mathematical checks, report, source-page inspections, exact-source offline build, and report-page inspections are current work. Detailed file hashes, original and recovery tokens, source/version bindings, limitations, and build/view records are retained with the per-paper verification files.\par
1 H. B. Bjerkevik, Stabilizing Decomposition of Multiparameter Persistence Modules, Foundations of Computational Mathematics 26 (2026), 939–998. Published online 27 January 2025. DOI: 10.1007/s10208-025-09695-w (https://doi.org/10.1007/s10208-025-09695-w).\par
\end{document}
